\chapter{La grille de lecture complète}
\label{chap:A-grille}

Cette annexe reproduit la grille de lecture du sondage, critère par critère, avec la question posée
au lecteur et sa règle de décision. Elle réunit la fiche du lecteur (\file{study/GRILLE.md}) et les
règles du protocole du sondage (\file{study/PROTOCOL.md}) telles qu'elles ont été appliquées,
harmonisation comprise, pour les grilles v1 et v2 et pour le critère complémentaire X8. En cas de
doute, le protocole fait foi.

\section{Ce que fait le lecteur}
\label{sec:a-lecteur}

Le lecteur lit l'article entier, annexes comprises. Il remplit un fichier JSON par article
(\cref{sec:formats}). Pour chaque critère, il donne la valeur, une citation exacte de l'article,
copiée mot pour mot du PDF, puis la page. Si rien ne justifie un \texttt{yes}, la citation est
\enquote{absent après lecture des sections expériences, annexes et liens} et la page vaut
\texttt{null}. Une citation sur deux pages reçoit une liste de pages. Dans un PDF à deux colonnes, une
citation qui chevauche des lignes se découpe en morceaux d'une ligne joints par \enquote{[...]}. Le
lecteur ne lit jamais les codes d'un autre lecteur, ni les résultats, ni le document de synthèse. Il
lance \file{study/check\_codes.py} sur ses fichiers et corrige jusqu'à zéro problème. Les valeurs
possibles sont \texttt{yes}, \texttt{no} et \texttt{na} (non applicable), sauf pour les critères
entiers et pour la famille X5.

\section{La grille de Pineau (Q1 à Q8)}
\label{sec:a-pineau}

Joelle Pineau a présenté ces huit critères à NeurIPS 2018 (diapositives 29 et 30, \cref{tab:a-pineau-2018}). Elle donnait pour
chacun la part des articles qui le satisfaisaient, sur 50 articles de RL de 2018 qu'elle avait lus. Sa méthode de tirage n'est
pas décrite. Ses intitulés sont donnés en anglais, tels qu'elle les a écrits. Le \cref{tab:a-grille-q}
donne la question posée et la règle de décision de chacun.

\begin{table}[htbp]
\centering
\small
\caption{Les huit critères de Pineau et ses chiffres de 2018.}
\label{tab:a-pineau-2018}
\begin{tabular}{@{}llr@{}}
\toprule
Code & Intitulé de Pineau & 2018 \\
\midrule
Q1 & Paper has experiments & 100~\% \\
Q2 & Paper uses neural networks & 90~\% \\
Q3 & All hyperparams for proposed algorithm are provided & 90~\% \\
Q4 & All hyperparams for baselines are provided & 60~\% \\
Q5 & Code is linked & 55~\% \\
Q6 & Method for choosing hyperparams is specified & 20~\% \\
Q7 & Evaluations on some variation of a hold-out test set & 10~\% \\
Q8 & Significance testing applied & 5~\% \\
\bottomrule
\end{tabular}
\end{table}

\begin{small}
\begin{longtable}{@{}p{0.07\linewidth}>{\raggedright\arraybackslash}p{0.27\linewidth}>{\raggedright\arraybackslash}p{0.6\linewidth}@{}}
\caption{Grille v1, critères de Pineau : question et règle de décision.}\label{tab:a-grille-q}\\
\toprule
Code & Question & Règle de décision \\
\midrule
\endfirsthead
\toprule
Code & Question & Règle de décision \\
\midrule
\endhead
\bottomrule
\endfoot
Q1 & L'article a-t-il des expériences ? & \texttt{no} : tous les autres critères valent \texttt{na}
  et sortent des dénominateurs. Les dénominateurs de Pineau sont les articles
  expérimentaux. \\
Q2 & Emploie-t-il des réseaux de neurones ? & Lecture directe. \\
Q3 & Tous les hyperparamètres de la méthode sont-ils donnés ? & \texttt{yes} si un tableau ou un
  paragraphe les donne pour toutes les méthodes principales. \enquote{Valeurs par défaut de
  [bibliothèque]} sans les lister compte \texttt{yes} seulement si la version est donnée. \\
Q4 & Tous les hyperparamètres des références sont-ils donnés ? & Même règle que Q3. \texttt{na}
  s'il n'y a aucune référence comparée. \\
Q5 & Le code est-il en lien ? & Un lien vers du code dans l'article. \enquote{Sur demande},
  \enquote{sera publié}, \enquote{will be made public at <URL>} : \texttt{no}. Une page de projet
  compte \texttt{yes} seulement si elle mène à du code, vérifié en la visitant ;
  un miroir anonyme qui contient le code compte \texttt{yes}. \\
Q6 & La façon de choisir les hyperparamètres est-elle décrite ? & La méthode de choix est dite :
  grille, recherche aléatoire, budget, critère. \enquote{Tuned} sans plus : \texttt{no}. \\
Q7 & L'évaluation se fait-elle sur des conditions tenues à part ? & Tâches, environnements, graines
  aléatoires d'environnement ou jeux de données non vus à l'entraînement ni au réglage. Pour un
  article \texttt{llm}, des bancs d'essai distincts des données d'entraînement RL comptent
  \texttt{yes}. Article \texttt{theory} à expériences jouets : \texttt{no}. Des graines aléatoires
  d'entraînement nouvelles, utilisées pour entraîner, ne sont pas des conditions tenues à part. \\
Q8 & Un test statistique est-il appliqué à une comparaison ? & Un test (t, Welch, Mann-Whitney,
  bootstrap avec intervalle, moyenne interquartile de rliable avec intervalles stratifiés) appliqué
  à une comparaison. Des barres d'erreur seules : \texttt{no} pour Q8, \texttt{yes} pour X2. \\
\end{longtable}
\end{small}

\section{Nos ajouts de la grille v1 (X1 à X7)}
\label{sec:a-x}

Ces critères viennent de la conférence d'ICLR 2018 et du débat de 2025 sur les récompenses trompeuses des
modèles de langue (\cref{sec:llm-pourquoi}). Ils font partie du protocole
(\cref{tab:a-grille-x}).

\begin{small}
\begin{longtable}{@{}p{0.07\linewidth}>{\raggedright\arraybackslash}p{0.27\linewidth}>{\raggedright\arraybackslash}p{0.6\linewidth}@{}}
\caption{Grille v1, nos ajouts : question et règle de décision.}\label{tab:a-grille-x}\\
\toprule
Code & Question & Règle de décision \\
\midrule
\endfirsthead
\toprule
Code & Question & Règle de décision \\
\midrule
\endhead
\bottomrule
\endfoot
X1 & Combien d'entraînements indépendants pour les résultats principaux ? & Un entier, ou
  \texttt{null} si non dit : le plus petit nombre annoncé pour un tableau ou une courbe principale.
  Pour un article \texttt{llm}, on compte les entraînements RL indépendants et non les échantillons
  d'évaluation : pass@k et \enquote{16 samples} ne sont pas des entraînements. \\
X2 & Une incertitude est-elle rapportée sur les résultats principaux ? & Écart type, erreur type,
  intervalle de confiance, bande autour de la courbe. Toute barre compte, quelle que soit la
  dispersion qu'elle résume : le critère X8 (\cref{sec:a-x8}) demande sur quoi elle est calculée. \\
X3 & Les meilleurs essais sont-ils sélectionnés ? & \enquote{Top-k}, \enquote{best of N}, meilleur
  point de contrôle choisi sur le test. \texttt{yes} aussi quand les propres tableaux de l'article
  montrent que le résultat principal est le meilleur de plusieurs entraînements identiques, même si
  le texte ne le dit pas. \\
X4 & L'article dit-il que les références sont réglées avec le même soin que la méthode ? & \texttt{yes} exige un
  réglage décrit pour les références, ou un budget de réglage égal annoncé. \enquote{Même
  configuration que la méthode}, sans réglage décrit pour personne : \texttt{no}. \\
X5 & À quelle famille appartient l'article ? & \texttt{llm} si la politique optimisée est un modèle
  génératif préentraîné conditionné par le langage (modèle de langue de toute taille, modèle de
  vision-langage, modèle texte-image). \texttt{classic} si le RL entraîne autre chose, même au service
  d'un grand modèle, ou un réseau transformeur entraîné de zéro sur des tâches de RL.
  \texttt{theory} pour un article sans expérience ou à expériences jouets seulement. Un article à
  résultats théoriques qui compare des méthodes sur des bancs d'essai usuels, même en annexe, est
  \texttt{classic}. \\
X6 & Pour \texttt{llm} : combien de familles de modèles de base distinctes ? & Un entier, ou
  \texttt{null}. Qwen, Llama, Mistral, Gemma, DeepSeek, OLMo\dots ; deux tailles de Qwen comptent pour
  une famille. \\
X7 & Le calcul employé est-il chiffré ? & \texttt{yes} exige au moins une quantité : heures,
  heures de GPU, FLOP, ou un nombre de GPU avec une durée. Le seul nom du matériel : \texttt{no}. \\
\end{longtable}
\end{small}

\section{La grille v2 (Y1 à Y12)}
\label{sec:a-y}

Les règles de décision de la grille v2 comprennent les précisions tirées du pilote de trois articles
et de l'audit des lectures brutes. Chaque critère renvoie à la source
de 2018 qui le motive.

\paragraph{Règles de portée} Si Q1 vaut \texttt{no}, Y1 à Y12 valent \texttt{na}. Y2 et Y11 valent
\texttt{na} si et seulement si X2 vaut \texttt{no}. Y4 vaut \texttt{na} si Y3 vaut 1. Y5 vaut
\texttt{na} si et seulement si Q4 vaut \texttt{na}. Pour un article déjà lu avec la grille v1, le
lecteur reçoit les valeurs v1 de Q1, Q4 et X2, qui fixent ces règles, sans voir les autres codes v1.
Le \cref{tab:a-grille-y} donne la question, la source et la règle de chaque critère Y.

\begin{small}
\begin{longtable}{@{}p{0.07\linewidth}>{\raggedright\arraybackslash}p{0.27\linewidth}>{\raggedright\arraybackslash}p{0.6\linewidth}@{}}
\caption{Grille v2 : question, source et règle de décision.}\label{tab:a-grille-y}\\
\toprule
Code & Question & Règle de décision \\
\midrule
\endfirsthead
\toprule
Code & Question & Règle de décision \\
\midrule
\endhead
\bottomrule
\endfoot
Y1 & La mesure de performance est-elle définie ? (ICLR 2018, diapositive 27) & \texttt{yes} si
  l'article dit la quantité et le point de mesure : retour final, moyenne sur les N derniers
  épisodes, meilleur point, aire sous la courbe, pass@k avec k et le nombre d'échantillons, taux de
  réussite et sur quoi. \enquote{Return} ou \enquote{accuracy} sans plus : \texttt{no}. Un score
  normalisé dont on ne dit pas à quel point il est pris : \texttt{no}. Le nombre d'épisodes
  d'évaluation n'est pas exigé. Une courbe dont les points d'évaluation sont dits vaut un
  point de mesure, \enquote{entraîné N pas} puis un seul chiffre vaut le point final. Pour une
  quantité qui n'est pas le score d'une politique entraînée (erreur d'un estimateur, temps de
  calcul), \texttt{yes} si la quantité est définie. \\
Y2 & Les barres d'erreur sont-elles définies ? (liste de contrôle v1.2) & \texttt{yes} si l'article
  dit ce que représentent la bande ou l'intervalle : écart type, erreur type, intervalle de confiance
  à 95~\%, minimum et maximum. Des barres présentes mais non définies : \texttt{no}. \\
Y3 & Combien de tâches la conclusion principale couvre-t-elle ? (Henderson et al., tableau 3) & Un
  entier : le nombre de cadres d'entraînement distincts sur lesquels la méthode est entraînée et
  comparée, un environnement ou un jeu de données d'entraînement pour \texttt{llm}. Deux niveaux de
  difficulté d'un même jeu, ou les variantes de qualité d'un jeu de données d'un même environnement
  (D4RL medium, replay, expert), comptent pour une tâche. Plusieurs bancs d'évaluation d'un même
  modèle entraîné ne multiplient pas les tâches. Un seul méta-entraînement compte pour une
  tâche. Un cadre présent seulement en annexe ne compte pas. Valeur dérivée pour
  l'analyse : \texttt{yes} si l'entier vaut au moins 2. \\
Y4 & Une tâche où la méthode ne gagne pas est-elle montrée ? (Henderson et al., tableau 3) &
  \texttt{yes} si un résultat tâche par tâche est donné et qu'au moins une tâche où la méthode ne
  gagne pas est montrée ou dite. \texttt{no} si l'article ne donne que des moyennes, ou ne montre que
  des tâches où la méthode gagne. \\
Y5 & L'implémentation des références est-elle identifiée ? (Henderson et al., figure 6) &
  \texttt{yes} si, pour chaque référence principale, l'article dit d'où vient le code : dépôt,
  bibliothèque ou cadre logiciel nommé, avec ou sans version, \enquote{réimplémenté par nous},
  chiffres recopiés d'articles avec leur source, modèle publié désigné par son nom exact, cadre
  d'entraînement des auteurs s'il est dit que les références y tournent. Une
  seule référence principale sans origine donne \texttt{no}. \\
Y6 & Les exclusions sont-elles dites ? (liste de contrôle v1.2) & \texttt{yes} si l'article dit
  quels essais, graines aléatoires ou données ont été retirés, ou dit explicitement qu'aucun ne l'a été. Une
  combinaison omise avec sa raison compte, un filtrage des données d'entraînement dit avec
  ce qu'il retire aussi. Le silence vaut \texttt{no}. \\
Y7 & La version des environnements ou des bancs d'essai est-elle donnée ? (liste de contrôle v1.2)
  & \texttt{yes} si la version de chaque environnement ou banc principal est donnée : suffixe Gym ou
  Gymnasium, version de la suite, nom exact et date de la séparation de test pour \texttt{llm}.
  \enquote{MuJoCo} seul : \texttt{no}. \enquote{MATH-500} compte \texttt{yes}, \enquote{MATH} seul
  \texttt{no}. Un jeu publié sur un dépôt nommé compte comme versionné, une séparation construite par
  les auteurs seulement si elle est publiée. Un nom qui couvre plusieurs ensembles sans
  dire lequel (\enquote{AMC}, \enquote{Minerva}) ne compte pas. Un environnement entièrement spécifié
  dans l'article compte. \\
Y8 & Un test de généralisation hors des conditions d'entraînement est-il fait ? (NeurIPS 2018,
  feuille de route) & \texttt{yes} si au moins une évaluation se fait sur des conditions absentes de
  l'entraînement et du réglage : autres graines aléatoires d'environnement annoncées comme telles, tâches
  nouvelles, dynamiques modifiées, distribution de test distincte. Plus strict que Q7 : pour
  \texttt{llm}, un banc différent ne compte que s'il est présenté comme un test hors distribution.
  Une séparation de test de même distribution relève de Q7. Y8 \texttt{yes} entraîne Q7 \texttt{yes}. \\
Y9 & La sensibilité à au moins un hyperparamètre est-elle montrée ? (ICLR 2018, diapositives 24 à
  26) & \texttt{yes} si au moins un hyperparamètre de la méthode, présenté comme tel, est balayé et
  que le résultat est montré par une mesure de performance. Une ablation de composant ne compte pas.
  Un paramètre réglé à l'inférence compte s'il est présenté comme un hyperparamètre. Une condition
  de l'environnement ne compte pas. \\
Y10 & Les dépendances logicielles sont-elles listées avec leurs versions ? (liste de contrôle v1.2)
  & \texttt{yes} si les versions des bibliothèques principales sont dans l'article, ou si un fichier
  d'environnement est cité dans le dépôt lié. Le dépôt n'est pas visité, sauf si l'article renvoie
  explicitement au fichier. \\
Y11 & Le type d'incertitude est-il nommé ? (comparable à Gorsane et al., 2022) & \texttt{yes} ou
  \texttt{no}, plus un champ \texttt{kind} : \texttt{std}, \texttt{se}, \texttt{ci}, \texttt{iqr},
  \texttt{minmax} ou \texttt{other}. Plus large que Y2 : il suffit que le type soit nommé quelque
  part. Des bandes de centiles autres que l'intervalle interquartile : \texttt{other}. \\
Y12 & Un outil statistique postérieur à 2018 est-il employé ? (adoption d'Agarwal et al., 2021) &
  \texttt{yes} si l'article emploie au moins un de ces outils : moyenne interquartile, intervalle
  bootstrap stratifié, rliable, profils de performance, probabilité d'amélioration. Un simple
  bootstrap sur une moyenne compte s'il produit un intervalle rapporté. \\
\end{longtable}
\end{small}

\section{Le critère complémentaire X8}
\label{sec:a-x8}

X2 compte toute barre d'erreur. Or une barre calculée sur les épisodes d'un même entraînement est le
piège que montre l'expérience B1 (\cref{sec:b1}). Le critère X8 (\cref{tab:a-grille-x8}) demande donc si l'incertitude est
calculée sur des entraînements indépendants. Il est exploratoire : il s'ajoute aux résultats sans
changer les verdicts de la règle à grille égale.

\begin{small}
\begin{longtable}{@{}p{0.07\linewidth}>{\raggedright\arraybackslash}p{0.27\linewidth}>{\raggedright\arraybackslash}p{0.6\linewidth}@{}}
\caption{Le critère X8 : question et règle de décision.}\label{tab:a-grille-x8}\\
\toprule
Code & Question & Règle de décision \\
\midrule
\endfirsthead
\toprule
Code & Question & Règle de décision \\
\midrule
\endhead
\bottomrule
\endfoot
X8 & L'incertitude des résultats principaux est-elle calculée sur des graines aléatoires ? &
  \texttt{yes} si l'incertitude rapportée sur un tableau ou une courbe principale est calculée entre
  au moins deux entraînements indépendants et que l'article le dit (\enquote{mean $\pm$ std over 5
  seeds}, \enquote{shaded area: standard error across 10 runs}). \texttt{no} si elle est calculée sur des épisodes, des déroulements d'évaluation ou des éléments de test seulement, si elle réunit dans un
  même écart les épisodes de toutes les graines aléatoires, ou si sa source n'est pas dite.
  \texttt{na} si X2 vaut \texttt{no}. Pour un article \texttt{llm}, on compte les entraînements RL
  indépendants, comme pour X1. Un champ obligatoire, \texttt{source}, nomme ce sur quoi la barre
  est calculée : \texttt{seeds}, \texttt{pooled} (graines aléatoires et épisodes réunis),
  \texttt{episodes}, \texttt{items}, \texttt{unstated} ou \texttt{other}. Citation exacte et page,
  comme pour tout critère. \\
\end{longtable}
\end{small}

X8 a été lu sur les 74 articles de la comparaison principale où X2 vaut \texttt{yes}, 30 de 2018 et
44 de 2026, par huit agents neufs, S1 à S8, qui n'avaient lu aucun article du sondage. Les six lecteurs
principaux, S1 à S6, ont reçu les articles tour à tour après un mélange \texttt{random.Random(2032)},
si bien que chacun a lu les deux années. S7 et S8 ont relu vingt articles tirés par
\texttt{random.Random(2033)}, dix par année. Les lecteurs recevaient la règle, pas l'hypothèse. Les
citations sont vérifiées par la fonction de \file{study/check\_codes.py}.

\section{L'harmonisation}
\label{sec:a-harmonisation}

Une règle d'harmonisation figure dans le protocole. Elle s'applique aux deux années, depuis les citations, dans \file{study/overrides.json}
(grille v1) ou \file{study/overrides\_v2.json} (grille v2), chaque valeur modifiée avec sa raison.
L'accord entre lecteurs se calcule sur les lectures brutes, avant harmonisation. Pour X3, X4 et X7,
aucune règle n'a été ajoutée après la lecture élargie, pour que les deux années restent lues avec la
même grille. La fonction \texttt{pineau\_rl.grid.validate} contrôle le schéma et les règles
de portée de chaque fichier (\cref{sec:formats}).
